\documentclass[a4paper,11pt,twoside,final]{book}

\usepackage[british]{babel}       % British English
\usepackage[T1]{fontenc}          % Required for hyphenation and \DJ
\usepackage{hyphenat}             % \hyp{} is a breakable dash
\usepackage[dvipsnames]{xcolor}
\usepackage{amsmath,amssymb,amsthm,stmaryrd}
\theoremstyle{plain}

\title{Sorting $2^n$ keys}
\author{Christian Rinderknecht}

\begin{document}

\maketitle

\paragraph{Minimum cost}

Let \(B(\Join,p)\) be the lowest cost of sorting \(2^p\) keys by merging
them. We have:
\begin{equation}
B(\Join, 0) = 0, \qquad B(\Join, p) = p2^{p-1}, \quad \text{for \(p > 0\)}.
\end{equation}

\paragraph{Maximum cost}

Let \(W(\Join,p)\) be the highest cost of sorting \(2^p\) keys by merging
them. We have:
\begin{equation}
W(\Join, p) = (p-1)2^p + 1, \quad \text{for \(p \geqslant 0\)}.
\end{equation}

\paragraph{Average cost}
\label{par:Atms_2p}

Let \(M(\Join, p)\) be the average cost of sorting \(2^p\) keys by merging
them. The trivial case is \(M(\Join,0) = 0\). We have, for \(p > 0\),
\begin{align*}
 M(\Join, p)
  &= 2^p\sum_{k=0}^{p-1}{\frac{1}{2^k}\left(2^k - 1 + \frac{1}{2^k +
      1}\right)}\\
  &= 2^{p}\left(p - \sum_{k=0}^{p-1}{\frac{1}{2^k}}
     + \sum_{k=0}^{p-1}{\frac{1}{2^k(2^k+1)}}\right)\\
  &= 2^{p}\left(p - \sum_{k=0}^{p-1}{\frac{1}{2^k}}
     + \sum_{k=0}^{p-1}\left(\frac{1}{2^k}
     - \frac{1}{2^k+1}\right)\!\!\right)
     = p2^p - 2^p \sum_{k=0}^{p-1}\frac{1}{2^k+1}.
\end{align*}
Let us work on the sum by splitting it and looking for a limit:
\begin{equation}
  \sum_{k=0}^{p-1}\frac{1}{2^k+1} = \sum_{k=0}^q\frac{1}{2^k+1} -
  \sum_{k=p}^q\frac{1}{2^k+1}, \quad \text{for \(q > p\)}.
  \label{eq:alpha_conv}
\end{equation}
This leads us to consider the convergence of \(\sum_{k \geqslant
  0}1/(2^k+1)\). We have the obvious bounds \(0 < 1/(2^k+1) < 2^{-k}\), so the
partial sums~\(\sum_{k=0}^q{1/(2^k+1)}\) are bounded by \(\sum_{k=0}^{q}2^{-k} <
2\). Moreover, those sums are strictly increasing, as they are made of positive
terms only, therefore \(\alpha := \sum_{k \geqslant 0}1/(2^k+1)\) exists and
\(\alpha \leqslant 2\). More precisely:
\begin{equation*}
  \alpha
  = \frac{1}{2^0+1} + \sum_{k \geqslant 1}\frac{1}{2^k+1}
  = \frac{1}{2} + \sum_{k \geqslant 1}\frac{1}{2^k+1}
  < \frac{1}{2} + \sum_{k \geqslant 1}\frac{1}{2^k}
  = \frac{1}{2} + 1 = \frac{3}{2}.
\end{equation*}
Borwein (1992) showed that \(\alpha \simeq 1.264500\) is irrational. Since the
convergence of \(\sum_{k=0}^q{1/(2^k+1)}\) implies that of
\(\sum_{k=p}^q{1/(2^k+1)}\), we can have \(q \rightarrow \infty\)
in~\eqref{eq:alpha_conv}:
\begin{equation*}
  \sum_{k=0}^{p-1}\frac{1}{2^k+1} = \sum_{k \geqslant 0}\frac{1}{2^k+1} -
  \sum_{k \geqslant p}\frac{1}{2^k+1}.
\end{equation*}
We can then resume our derivation of \(M(\Join, p)\):
\begin{align}
 M(\Join, p)
  &= p2^p - 2^p \sum_{k \geqslant 0}\frac{1}{2^k+1}
   + 2^p \sum_{k \geqslant p}\frac{1}{2^k+1}\notag\\
  &= p2^p - 2^p \alpha + \sum_{k \geqslant p}\frac{1}{2^{k-p}+2^{-p}}\notag\\
  &= p2^p - 2^p \alpha + \sum_{k \geqslant 0}\frac{1}{2^k+2^{-p}}
   = (p - \alpha) 2^p + r_p,
\label{eq:Mjoin}
\end{align}
where \(r_p := \sum_{k \geqslant 0}1/(2^{k}+2^{-p})\). We have:
\begin{equation*}
  \frac{1}{2^k + 2^{-p-1}} - \frac{1}{2^k + 2^{-p}}
  = \frac{2^{-p-1}}{(2^k + 2^{-p-1})(2^k + 2^{-p})} > 0,
\end{equation*}
so we get, by summation,
\begin{equation*}
  r_{p+1} - r_p = 2^{-p-1}\sum_{k \geqslant 0}\frac{1}{(2^k + 2^{-p-1})(2^k +
    2^{-p})} > 0,
\end{equation*}
so \(r_p\)~is strictly increasing. Let us now show that \(2\)~is its limit
and a strict upper bound, by examining the difference:
\begin{equation*}
  2 - r_p
  = \sum_{k \geqslant 0}\left(\frac{1}{2^k} - \frac{1}{2^k + 2^{-p}}\right)
  = 2^{-p}\sum_{k \geqslant 0}\frac{1}{2^k(2^k + 2^{-p})}
  \leqslant 2^{-p}\sum_{k \geqslant 0}\frac{1}{4^k}
  = \frac{4}{3}2^{-p}.
\end{equation*}
The series on the left is made of positive terms, so \(0 < 2 - r_p \leqslant
\tfrac{4}{3}2^{-p}\) and, since \(\tfrac{4}{3}2^{-p} \rightarrow 0\) as \(p \to
\infty\), we deduce \(r_p \rightarrow 2^{-}\), thus \(M(\Join, p) < (p - \alpha)
2^p + 2\) for \(p \geqslant 0\). This bound is optimal in the sense that the
constant~\(2\) cannot be lowered since \(r_p \rightarrow 2^{-}\). Also
\(M(\Join,p)\) is best approximated by it for sufficiently large values
of~\(p\): \(M(\Join, p) = (p - \alpha) 2^p + 2 - \theta_p\), with \(0 < \theta_p
\leqslant \frac{4}{3}2^{-p}\). Now, let us turn our attention to lower
bounds. We established that \(r_p\) is strictly increasing for \(p \geqslant
0\), but since we leave aside the trivial case \(p=0\) (sorting one key requires
no comparison), the greatest lower bound is~\(r_1 = \min_{p \geqslant 1}r_p\):
\begin{equation*}
  r_1 = \sum_{k \geqslant 0}\frac{1}{2^k + 2^{-1}}
  = \sum_{k \geqslant 0}\frac{2}{2^{k+1} + 1}
  = 2 \sum_{m \geqslant 1}\frac{1}{2^m + 1}
  = 2\left(\!\alpha - \frac{1}{2^0 + 1}\!\right)
  = 2\alpha - 1.
\end{equation*}
We conclude
\begin{equation}
(p - \alpha)2^p + 2\alpha - 1 \leqslant M(\Join,p) < (p-\alpha)2^p + 2, \quad
  \text{for \(p > 0\),}
\end{equation}
where the lower bound is an equality if, and only if, \(p = 1\), and the upper
bound holds even if \(p = 0\).
\end{document}
